\documentclass[10pt, letterpaper]{article}

% --- Packages ---
\usepackage[
    ignoreheadfoot,
    top=1.0 cm,
    bottom=1.0 cm,
    left=1.5 cm,
    right=1.5 cm,
    footskip=1.0 cm,
]{geometry}
\usepackage{titlesec}
\usepackage{fontawesome5}
\usepackage[dvipsnames]{xcolor}

% --- Mélissa's signature red ---
\definecolor{MyRed}{RGB}{160, 20, 40}

\usepackage[
    pdftitle={Melissa Colin - Cover Letter - AMI Labs},
    pdfauthor={Melissa Colin},
    colorlinks=true,
    urlcolor=MyRed
]{hyperref}
\usepackage{charter}

% --- Settings ---
\pagestyle{empty}
\setlength{\parindent}{0pt}
\setlength{\parskip}{0.8em}

\IfFileExists{personal.tex}{\input{personal.tex}}{}
\providecommand{\myphone}{+33 7 82 52 XX XX}

\begin{document}

% ==================== HEADER ====================
\begin{center}
    {\fontsize{22}{22}\selectfont \textbf{\color{MyRed}Mélissa Colin}} \\[4pt]
    {\large \textbf{AI Research, Physics-grounded Motion Generation, World Models}} \\[4pt]
    \small{
    \faMapMarker* Bordeaux \quad|\quad
    \faPhone* \myphone{} \quad|\quad
    \faEnvelope~ \href{mailto:contact@melissacolin.ai}{Email} \quad|\quad
    \faGlobe~ \href{https://melissacolin.ai}{Portfolio} \quad|\quad
    \faLinkedin~ \href{https://www.linkedin.com/in/melissacolin/}{LinkedIn} \quad|\quad
    \faGithub~ \href{https://github.com/melissa-colin}{GitHub}
    }
\end{center}

\vspace{12pt}

% ==================== RECIPIENT ====================
\textbf{AMI Labs} \\
General Application 2026--2027 \\
Paris, France

\textbf{Subject: Spontaneous application for a research position (apprenticeship Sept 2026, or internship Feb 2027)}

% ==================== BODY ====================
Dear AMI Labs Team,

I want to do research on AI architectures that work closer to how biological intelligence does. AMI Labs is the only place I've found where that's the central agenda and not a side topic. The JEPA and world models direction is what made me apply.

I'm currently a Visiting Research Scholar at the University of Alberta Vision and Learning Lab (Mitacs Globalink Award 2026), with Prof.~Li Cheng and mentored by Yilin Wang (first author of MotionDreamer, ICLR 2025). I work on physics-grounded multi-person motion generation: making generated human motion satisfy real physical and biomechanical constraints, not just look right on a screen. I'm working with Yilin as a co-author on a paper in preparation that we're aiming to submit to CVPR 2027.

Two earlier experiences pushed me toward architecture-level research. A summer R\&D internship at Sector Group on industrial AI safety, where I built an offline on-premise RAG pipeline (Ollama + BGE embeddings, 70\,\% recall, 100\,\% reliability) and worked on the IEC 61508 functional-safety standard. And a first scientific paper I wrote on my free time with Ikram Chraibi Kaadoud to learn how research works (published at PFIA 2024); we compared LIME, SHAP and Grad-CAM and found the three disagreed on the same models.

On the engineering side, 14 months at Cali Intelligences on YOLO-based detection: I designed and built the company's MLOps training pipeline from scratch, brought the production model from 60\,\% to 90\,\% accuracy on a 100,000-video dataset, and cut inference latency by 70\,\%.

I'd like to join AMI as either a 6-month research intern starting February 2027 (end-of-studies internship from ENSEIRB-MATMECA's Engineer-Doctor research track, an engineering degree with a research focus that prepares for a PhD; Top 3\,\% rank), or on a 12-month French work-study contract (contrat de professionnalisation) starting September 2026. The role that would fit me best is one bridging research and engineering.

Thank you for considering my application.

\vspace{6pt}
Sincerely, \\
\textbf{Mélissa Colin}

\end{document}
